% =====================================================================
%  Section III. Lattice Quantization of Gauge Fields
% =====================================================================
\section{LATTICE QUANTIZATION OF GAUGE FIELDS}\label{sec:lattice}

\subsection{Classical action on a lattice}\label{sec:classical}

In this section the gauge-field action (space-time integral of the
Lagrangian) will be defined on a discrete lattice with spacing $a$ in
Euclidean space-time.  The simplest way to proceed is to consider a
continuum action, substitute finite-difference approximations for
derivatives, and replace the space-time integral by a sum over the
lattice sites.  However, the result of this is an action which is not
gauge-invariant for nonzero $a$.  Because of the vagaries of
renormalization this is likely to mean that the quantized theory still
lacks gauge invariance in the limit $a \to 0$ (if such a limit exists).
The alternative is to formulate gauge invariance for a lattice theory,
and tinker with the action so that it is gauge invariant for any $a$.
This alternative will be pursued here.

For convenience a charged Dirac field $\psi$ coupled to a single gauge
field $A_{\mu}$ will be discussed in detail.  Generalizations to
non-Abelian gauge groups will be noted later.

On the lattice the fields are $\psi_{n}$, $\bar\psi_{n}$, and
$\thl{\mu\nu}(n)$, where $n$ is a four-vector with integer components
referring to points on a simple hypercubic lattice.

A simple action on the lattice for the Dirac field (ignoring the gauge
field for now) has the form
\begin{equation}
A_{D} = -a^{3}\sum_{n} m_{b}\bar\psi_{n}\psi_{n}
        + a^{3}\sum_{n,\mu} \bar\psi_{n}\gamma_{\mu}
          \frac{\psi_{n+\hat\mu}-\psi_{n-\hat\mu}}{2a},
\label{eq:3.1}
\end{equation}
where $m_{b}$ is the bare mass; $\hat\mu$ is a unit vector along the axis
$\mu$; $a^{3}\sum_{n}$ replaces the space-time integration of the
continuum theory, and $(\psi_{n+\hat\mu}-\psi_{n-\hat\mu})/2a$ replaces
$\partial_{\mu}\psi$.  There is no over-all factor of $i$ due to the
Euclidean metric.  A gauge transformation on the lattice can be defined
as follows.
\begin{equation}
\psi_{n} \to e^{i\Lambda_{n}}\psi_{n},\qquad
\bar\psi_{n} \to \bar\psi_{n}e^{-i\Lambda_{n}},
\label{eq:3.2}
\end{equation}
\begin{equation}
\thl{\mu\nu}(n) \to \thl{\mu\nu}(n) - \frac{\gamma_{n+\hat\mu}-\gamma_{n}}{a},
\label{eq:3.3}
\end{equation}
\stepcounter{equation}%   % the paper has no Eq. (3.4)
where $g$ is the coupling constant and $\gamma_{n}$ is arbitrary.  The
terms $\bar\psi_{n}\psi_{n+\hat\mu}$ and $\bar\psi_{n+\hat\mu}\psi_{n}$
are not invariant to this transformation; the corresponding
gauge-invariant expressions are
$\bar\psi_{n}\psi_{n+\hat\mu}\exp(iga\,\thl{\mu\nu}(n))$ and
$\bar\psi_{n}\psi_{n+\hat\mu}\exp(-iga\,\thl{\mu\nu}(n))$.  Thus, a
gauge-invariant form for $A_{D}$ is
\begin{equation}
\begin{split}
A_{D} = {} & -a^{3}\sum_{n} m_{b}\bar\psi_{n}\psi_{n}\\
           & + \frac{a^{2}}{2}\sum_{n,\mu}
             \Bigl[\bar\psi_{n}\gamma_{\mu}\psi_{n+\hat\mu}
             e^{iga\,\thl{\mu\nu}(n)}
             -\bar\psi_{n+\hat\mu}\gamma_{\mu}\psi_{n}
             e^{-iga\,\thl{\mu\nu}(n)}\Bigr].
\end{split}
\label{eq:3.5}
\end{equation}

In the action $A_{D}$, the field $\thl{\mu\nu}$ acts like an angular
variable:  $A_{D}$ is periodic in $\thl{\mu\nu}$ with period $2\pi$.

The free gauge-field action will be defined to preserve this property.
This does not mean that $\thl{\mu\nu}$ is an angular variable in the
continuum limit.  Owing to the relation~\eqref{eq:3.6},
$\thl{\mu\nu}$ becomes an infinitesimal angular variable $A_{\mu}(x)$
for $a \to 0$; such a variable has the range
$-\infty < A_{\mu}(x) < \infty$ without any periodicity.
\begin{equation}
\thl{\mu\nu}(n) = a\,g\,A_{\mu}(na) + \order{a^{2}}.
\label{eq:3.6}
\end{equation}

A gauge-invariant lattice approximation for $F_{\mu\nu}$ is
\begin{equation}
F_{\mu\nu}(n) =
\frac{\thl{\mu\nu}(n)-\thl{\nu\mu}(n)-\thl{\mu\nu}(n+\hat\nu)
+\thl{\nu\mu}(n+\hat\mu)}{a}.
\label{eq:3.7}
\end{equation}

It is convenient to define a rescaled form of $\thl{\mu\nu}$:
\begin{equation}
f_{\mu\nu}(n) = \thl{\mu\nu}(n)+\thl{\nu\mu}(n+\hat\mu)
-\thl{\mu\nu}(n+\hat\nu)-\thl{\nu\mu}(n).
\label{eq:3.8}
\end{equation}

A simple lattice action for the gauge field which preserves periodicity
is
\begin{equation}
A_{G} = \frac{1}{2g^{2}}\sum_{n,\mu\nu}
        \Bigl[1-\tfrac{1}{2}\bigl(e^{if_{\mu\nu}(n)}+e^{-if_{\mu\nu}(n)}\bigr)\Bigr].
\label{eq:3.9}
\end{equation}

In the continuum limit, $f_{\mu\nu} \to 0$ due to the factor $a$ in the
definition of $f_{\mu\nu}$.  Thus for small $a$, one can write
\[
A_{G} = \frac{1}{2g^{2}}\sum_{n,\mu\nu}
        \Bigl[\tfrac{1}{4}\bigl(f_{\mu\nu}(n)\bigr)^{2}+\cdots\Bigr],
\]
The constant term is irrelevant.  The linear term in $f_{\mu\nu}$ is 0
because $f_{\mu\nu}$ is odd in the indices $\mu$ and $\nu$.  The
quadratic term gives
\stepcounter{equation}%   % the paper has no Eq. (3.10)
\begin{equation}
A_{G} = -\frac{1}{2g^{2}}\,a^{4}\sum_{n,\mu\nu}\bigl(F_{\mu\nu}(n)\bigr)^{2},
\label{eq:3.11}
\end{equation}
which is the conventional gauge-field action in a lattice approximation.
The terms involving $f_{\mu\nu}^{3}$, $f_{\mu\nu}^{4}$, etc.\ all vanish
for $a\to0$ even after removing a factor $a^{4}$ to convert
$\sum_{n}$ into an integral.

The full action may now be written
\begin{equation}
\begin{split}
A = {} & -\epsilon\sum_{n}\bar\psi_{n}\psi_{n}
         + \kappa\sum_{n,\mu}\Bigl[
           \bar\psi_{n+\hat\mu}\gamma_{\mu}\psi_{n}e^{i\thl{\mu\nu}(n)}
           -\bar\psi_{n}\gamma_{\mu}\psi_{n+\hat\mu}e^{-i\thl{\mu\nu}(n)}\Bigr]\\
       & + \frac{1}{2g^{2}}\sum_{n,\mu\nu}
         \Bigl[1-\tfrac{1}{2}\bigl(e^{if_{\mu\nu}(n)}+e^{-if_{\mu\nu}(n)}\bigr)\Bigr],
\end{split}
\label{eq:3.12}
\end{equation}
with $\epsilon = m_{0}a^{4}$, $\kappa = a^{2}/2$.  This action reduces
to the usual continuum action for $a \to 0$; for finite $a$ it is gauge
invariant and periodic in the gauge field.  Note, however, that the
continuum limit is a classical limit in which the lattice variables
$\psi_{n}$, $\bar\psi_{n}$, and $\thl{\mu\nu}$ approach continuum
functions $\psi(x)$, $\bar\psi(x)$, and $A_{\mu}(x)$ with $x=na$.  The
continuum limit of the quantized theory is much harder to discuss owing
to renormalization problems.

\subsection{Quantization}\label{sec:quantization}

The problem of principal interest here is the quantization of the gauge
field.  Therefore, the gauge field will be quantized by itself to start
with.  Later the quantization of $\psi$ will be discussed.  At the end of
this section the generalizations to non-Abelian gauge theories will also
be described.

The quantization of the lattice gauge theory will be carried out in two
steps.  The first step will be to define a lattice version of Euclidean
vacuum expectation values, starting from a lattice version of the Feynman
path integral.  The second step will be to define a quantum theory on the
lattice, which will allow the introduction of a real time variable and the
definition of particle states and scattering amplitudes.  In both cases
the lattice provides an ultraviolet cutoff and there is no Lorentz
invariance.  Lorentz invariance can only be achieved in the limit
$a$ (lattice spacing) $\to 0$, if at all, and in practice this is a
difficult limit to evaluate.

As discussed in Sec.~\ref{sec:intro}, one would like to calculate the
gauge-field average of $\exp[ig\oint A_{\mu}(x)\,ds^{\mu}]$ weighted
with the gauge-field action.  On a lattice the line integral becomes a
sum over a closed path $P$ on the lattice (see, e.g.,
Fig.~\ref{fig:8}).  The sum has the form $ig\sum_{P}(\pm)\thl{\mu\nu}$,
where a particular $\thl{\mu\nu}$ is present in the sum if the path
connects the sites $n$ and $n+\hat\mu$ ($-\thl{\mu\nu}$ appears if the
path goes from $n+\hat\mu$ to $n$).

\begin{figure}[t]
  \centering
  \includegraphics[width=0.4\textwidth]{fig8}
  \refstepcounter{figure}%
  \caption*{Fig. 8. Example of a lattice path $P$.}
  \label{fig:8}
\end{figure}

On the lattice, an average over all gauge fields involves integrating
over all values of the $\thl{\mu\nu}$ for all $n$ and $\mu$.  Normally
one would have integrals over an infinite range:
$-\infty < \thl{\mu\nu} < \infty$, but because of the periodicity in
$\thl{\mu\nu}$ there is no point to integrating over more than a single
period.  Thus the lattice version of the gauge-field average is
\begin{equation}
I_{\theta}(P) = \prod_{n,\mu\nu}\int\ddth{\mu\nu}
  \exp\!\Bigl[i\sum_{P}(\pm)\thl{\mu\nu}
    +\frac{1}{2g^{2}}\sum_{n,\mu\nu}
      \bigl(e^{if_{\mu\nu}(n)}+e^{-if_{\mu\nu}(n)}\bigr)\Bigr].
\label{eq:3.13}
\end{equation}
\begin{equation}
Z_{\theta} = \prod_{n,\mu\nu}\int\ddth{\mu\nu}
  \exp\!\Bigl[\frac{1}{2g^{2}}\sum_{n,\mu\nu}
    \bigl(e^{if_{\mu\nu}(n)}+e^{-if_{\mu\nu}(n)}\bigr)\Bigr].
\label{eq:3.14}
\end{equation}

Note that no gauge-fixing term has been added to the action.  The finite
range of $\thl{\mu\nu}$ makes a gauge-fixing term unnecessary.  In
continuum gauge theories where $A_{\mu}(x)$ has an infinite range
$[-\infty < A_{\mu}(x) < \infty]$ a gauge-fixing term is essential to
have a convergent functional integral.  The reason for this is that the
volume in path-integral space generated by all possible gauge
transformations is infinite; the gauge-fixing term provides a convergence
factor in this volume.  In the lattice theory the total volume of
integration is finite if the lattice itself is of finite extent; no
convergence factor is required.  For a lattice of infinite extent there
are divergences due to the infinite number of integrations (in other
words, the infinite lattice volume), but these divergences are normally
removed by the division by $Z$ (this division is equivalent to removing
all vacuum loops in perturbation theory).

One can define more conventional vacuum expectation values in a similar
manner.  For instance, one can define a propagator.  In the absence of a
gauge-fixing term it is awkward to define a propagator for the gauge
field $\thl{\mu\nu}$ itself; instead one can define a gauge-invariant
propagator as
\begin{equation}
D_{\mu\nu}(n) = Z_{\theta}^{-1}\prod_{m,\rho}\int\ddth{\rho\sigma}
  e^{i\thl{\mu\nu}(n)}e^{-i\thl{\mu\nu}(0)}
  \exp\!\Big[\tfrac{1}{2g^{2}}\sum_{m,\rho\sigma}
    \bigl(e^{if_{\rho\sigma}(m)}+e^{-if_{\rho\sigma}(m)}\bigr)\Big].
\label{eq:3.15}
\end{equation}
This is a propagator for the operator $e^{i\thl{\mu\nu}}$; it is defined
only for the lattice points $n$ of a Euclidean space-time lattice.  If
the lattice spacing is $a$, this means the propagator is defined only for
imaginary times of the form $in_{0}a$, where $n_{0}$ is an integer.

A theory defined only for discrete imaginary values of the time leaves
much to be desired.  Fortunately, one can generalize the theory to define
a Hamiltonian for a quantized theory.  The particle eigenstates and
scattering amplitudes of the theory can then be obtained, in principle,
by diagonalizing the Hamiltonian.  The Hamiltonian will be defined using
the transfer-matrix formalism.  Only a brief discussion of the
transfer-matrix approach will be given here.  For a review${}^{6}$ of the
ideas see Wilson and Kogut.${}^{6}$  A detailed discussion including
approximate calculation of single-particle energies and scattering
amplitudes in a simple scalar field theory is given by Suri.${}^{5}$

Consider the expression for $Z$.  Introduce finite bounds on the lattice
coordinate $n_{0}$, say
\begin{equation}
-N \le n_{0} \le N.
\label{eq:3.16}
\end{equation}
Introduce periodic boundary conditions (see below${}^{6}$).  Then one can
write $Z$ as the trace of a matrix $V$, more precisely
\begin{equation}
Z = \Tr V^{2N+1}.
\label{eq:3.17}
\end{equation}
This formula is made possible by the fact that each term in the action
$A$ involves no more than two adjacent values of $n_{0}$.

To set up the matrix $V$, one must first understand the space on which it
acts.  The space used here is the space of all functions
$\Phi(\thl{ij})$ (periodic in each $\thl{ij}$ with period $2\pi$), where
the index $i$ runs from 1 to 3 only, and the lattice variable $n$ has
only three components $(n_{1},n_{2},n_{3})$.  The matrix $V$ will be
defined as a function of two sets of arguments, say $\thl{ij}$ and
$\thl{ij}'$, these two sets of arguments referring to the space-time
fields $\thl{ij}$ for two adjacent values of $n_{0}$.  Matrix
multiplication of two $V$'s involves integrations over a set of variables
$\thl{ij}$.  Define $V$ as
\begin{equation}
V = \prod_{n}\prod_{i,j}\int\ddth{i0}\,e^{U}.
\label{eq:3.18}
\end{equation}

The quantity $U$ as written out below looks complicated, but all it is is
that part of the action $A$ referring to a given nearest-neighbor pair of
values for $n_{0}$; terms in $A$ referring to a single value of $n_{0}$
are divided equally between the matrices connecting $n_{0}$ to
$n_{0}+1$ and $n_{0}$ to $n_{0}-1$.  The result is
\begin{equation}
U = \frac{1}{2g^{2}}\sum_{n,i,j}\Bigl[
    \bigl(e^{if_{0i,j}(n)}+e^{-if_{0i,j}(n)}\bigr)
    +\bigl(e^{i\thl{i0,j}(n)}+e^{-i\thl{i0,j}(n)}\bigr)\Bigr],
\label{eq:3.19}
\end{equation}
where
\begin{equation}
f_{0i}(n) = \thl{0i}(n)+\thl{i0}(n+\hat0)-\thl{0i}(n+\hat i)-\thl{i0}(n),
\label{eq:3.20}
\end{equation}
and
\begin{equation}
f_{i0}(n) = \thl{i0}(n)+\thl{0i}(n+\hat i)-\thl{i0}(n+\hat0)-\thl{0i}(n).
\label{eq:3.21}
\end{equation}

With the definition of $V$ given here, the trace $\Tr V^{2N+1}$ is easily
seen to reproduce all the integrations involved in the equation for $Z$,
and the sum of the $2N+1$ exponents $U$ reproduces the action $A$ except
for some additional terms coupling the boundary $n_{0}=N$ to the
boundary $n_{0}=-N$ to achieve a periodic structure.

Note that the 0 components of $\thl{\mu\nu}$ have received special
treatment.  This is because there are no terms in $A$ involving
$\thl{0\mu}$ for more than one value of $n_{0}$; this makes it possible
to include integrations over $\thl{0\mu}$ in the definition of $V$
rather than in the definition of matrix multiplication.

The matrix $V$ is used to define the quantized theory.  Briefly, this is
accomplished as follows.  $V$ is a Hermitian matrix, i.e.,
\begin{equation}
V(\theta,\theta')^{*} = V(\theta',\theta).
\label{eq:3.22}
\end{equation}
[This result can be verified by close examination of
Eqs.~\eqref{eq:3.18}--\eqref{eq:3.21}.  One must remember that the
variables $\thl{0\mu}$ are integrated out in the definition of $V$.  This
means that in forming the complex conjugate of $V$ one can also make the
change of variable $\thl{0\mu} \to -\thl{0\mu}$.]  Hence $V$ has a
complete orthogonal set of eigenstates $\psi_{\lambda}$ and eigenvalues
$\lambda$.  The Hamiltonian $H$ is now defined as follows:  The
eigenstates of $H$ are the eigenstates of $V$; the eigenvalues of $H$ are
given by
\begin{equation}
E_{\lambda} = -a^{-1}\ln\lambda,
\label{eq:3.23}
\end{equation}
where $\lambda$ is the corresponding eigenvalue of $V$.  The reason for
the factor $a^{-1}$ will be evident shortly.  The reason for using the
logarithm is so that the energies of multiparticle scattering states will
be the sum of single-particle energies (see Suri${}^{5}$).

A problem arises with this definition if $V$ has any negative eigenvalues
$\lambda$.  If this were to happen, $H$ would have complex eigenvalues.
This did not happen in the case studied by Suri${}^{5}$; whether it
happens here the author does not know; this question must be studied
further.  Even if $V$ has negative eigenvalues, they may be irrelevant in
the limit $a\to0$ if such a limit exists.

The definition~\eqref{eq:3.23} means that
\begin{equation}
V = e^{-aH}.
\label{eq:3.24}
\end{equation}
This means $V$ is the operator which propagates a state through an
imaginary time $ia$.  It is a consequence of this that the propagator
$D_{\mu\nu}$ is a vacuum expectation value for imaginary time $in_{0}a$
in the theory with Hamiltonian $H$.  For proof of this statement see
Refs.~5 and~6.

The lattice quantization procedure can be extended to non-Abelian gauge
theories.  This is done as follows.  In place of the single variable
$\thl{\mu\nu}$, one has a set of variables $\thl{\mu\nu}^{a}$ where $a$
is an internal index.  For each $n$ and $\mu$, $\thl{\mu\nu}$ is to
parametrize an element $h_{\mu\nu}$ of the gauge group.  In place of
$\exp(i\thl{\mu\nu})$ one substitutes $U(\thl{\mu\nu})$, where $U$ is the
unitary matrix representing $\thl{\mu\nu}$ in the quark representation.
The product $\bar\psi_{n}\gamma_{\mu}\psi_{n+\hat\mu}\exp(i\thl{\mu\nu})$
is replaced by $\bar\psi_{n}\gamma_{\mu}U(\thl{\mu\nu})\psi_{n+\hat\mu}$.
A gauge transformation is defined by a set of group transformations
$\gamma_{n}$; under these transformations
\begin{equation}
\thl{\mu\nu} \to \thl{\mu\nu}^{a},
\label{eq:3.25}
\end{equation}
\begin{equation}
h_{\mu\nu} = \exp\!\Bigl(i\sum_{a}T^{a}\thl{\mu\nu}^{a}\Bigr),
\label{eq:3.26}
\end{equation}
\begin{equation}
\psi_{n} \to U(\gamma_{n})\psi_{n},\qquad
\bar\psi_{n} \to \bar\psi_{n}U(\gamma_{n})^{-1},
\label{eq:3.27}
\end{equation}
\begin{equation}
U(\thl{\mu\nu}) \to U(\gamma_{n})\,U(\thl{\mu\nu})\,U(\gamma_{n+\hat\mu})^{-1},
\label{eq:3.28}
\end{equation}
where $\gamma_{n}\,\thl{\mu\nu}\,\gamma_{n+\hat\mu}^{-1}$ is computed
according to the multiplication law of the group.  A simple
gauge-invariant action for the gauge field is
\begin{equation}
\begin{split}
A_{G} = \frac{1}{2g^{2}}\sum_{n,\mu\nu}\Bigl[1-\tfrac{1}{2}\Tr
    &\bigl(U(\thl{\mu\nu}(n))\,U(\thl{\nu\mu}(n+\hat\mu))\\
    &\qquad U^{-1}(\thl{\mu\nu}(n+\hat\nu))\,U^{-1}(\thl{\nu\mu}(n))\bigr)\Bigr],
\end{split}
\label{eq:3.29}
\end{equation}
where the unitary transformation $U_{A}$ is taken from the adjoint
representation of the group.  (Any representation will do as well.)  In
the classical continuum limit, $\thl{\mu\nu}$ becomes an infinitesimal
group transformation $\thl{\mu\nu} = ga\,A_{\mu}^{a}(na)T^{a}$ with
$A_{\mu}^{a}$ fixed as $a\to0$; in this limit one can show, with some
effort, that $A_{G}$ reduces to the standard continuum Yang-Mills action.
When the non-Abelian theory is quantized, the integrations over
$\thl{\mu\nu}$ are group integrations over all elements of the group.
These are compact integrations (for compact groups) and no gauge-fixing
term is required.

Since quark binding can be illustrated using Abelian theory, the
non-Abelian theory will not be studied further.

Finally, the quantization of the Dirac field will be discussed.  It is
convenient initially to quantize the Dirac field in an analogous manner
to the gauge field.  The only problem that occurs is to define
``integration'' for a Fermi field.  This can be done.${}^{11}$

The property of the integral that is crucial for quantization is
translational invariance in the integration variable.  For example, when
quantizing a scalar field $\phi$ on a lattice the field averaging
involves the integral $\int\dd\phi_{n}$ which has the translational
invariance
\begin{equation}
\int\dd\phi_{n}\,f(\phi_{n}+\xi_{n}) = \int\dd\phi_{n}\,f(\phi_{n})
\label{eq:3.30}
\end{equation}
for any integrable function $f$ and any constant $\xi_{n}$.  It is this
translational invariance that makes the Feynman path integral provide a
realization of Schwinger's action principle (see, e.g., Ref.~12 for
further discussion).  Analogously one needs to define an integration over
Fermi fields with the same translational invariance.  Stated abstractly,
one wants to define a bracket operation $(\ )$ defined on functions of
purely anticommuting Fermi fields $\psi_{n}$ with the property
\begin{equation}
(f) = \int f(\psi,\bar\psi)\,\prod_{n}\dd\psi_{n}\,\dd\bar\psi_{n},
\label{eq:3.31}
\end{equation}
where $\psi_{n}$ and $\bar\psi_{n}$ are anticommuting $c$-numbers (these
have been introduced by Schwinger).  The bracket operation should produce
a number for every function $f$; it should also be a linear operation.
Thus for a finite lattice it is sufficient to specify the bracket
$(\ )$ for all monomials in the $\psi_{n}$ and $\bar\psi_{n}$.  Because of
the anticommutation rules, $\psi_{n}^{2}$ and $\bar\psi_{n}^{2}$ vanish
(more correctly, $\psi_{n\alpha}^{2}$ and $\bar\psi_{n\alpha}^{2}$ vanish
where $\psi_{n\alpha}$ and $\bar\psi_{n\alpha}$ are any component of
$\psi_{n}$ and $\bar\psi_{n}$), therefore, there are only a finite number
of possible monomials.  It is now easy to see that the bracket $(\ )$
must vanish for all products of $\psi$'s and $\bar\psi$'s, except the
product containing all possible $\psi$'s and $\bar\psi$'s.  For example,
suppose there are two lattice sites $0$ and $1$, and $\psi_{n}$ and
$\bar\psi_{n}$ are single component fields.  Then the brackets must be
\begin{equation}
(\psi_{0}\bar\psi_{0}\psi_{1}\bar\psi_{1}) = 1,
\label{eq:3.32}
\end{equation}
where 1 is a constant which was chosen arbitrarily.  This definition of
the bracket operation satisfies translational invariance:  for example,
because the terms multiplying the $\psi$'s are all 0.  Note also that
the anticommutation rules mean that, for example,
\begin{equation}
(\bar\psi_{0}\psi_{0}\bar\psi_{1}\psi_{1}) = -1.
\label{eq:3.33}
\end{equation}
(In analogy to the scalar case, one requires $\psi\bar\psi = 1$, not
$\psi\bar\psi = 1 - \bar\psi\psi$.)

One can now define the Feynman path integral on a lattice for the
complete gauge theory including the Dirac fields.  For example, the
current-current propagator on the lattice is
\begin{equation}
D_{\mu\nu}(x) = Z^{-1}\prod_{n,\rho}\int\dd\psi_{n}\,\dd\bar\psi_{n}\,
  \dd\thl{\rho}\,J_{\mu}(x)J_{\nu}(0)\,e^{A},
\label{eq:3.34}
\end{equation}
where $A$ is the full action of Eq.~\eqref{eq:3.12} and
\begin{equation}
J_{\mu}(n) = \bar\psi_{n}\gamma_{\mu}\psi_{n}.
\label{eq:3.35}
\end{equation}

This formulation of the path integral is different from the formulation
discussed in Sec.~\ref{sec:binding}.  However, one can easily derive a
lattice form of the path integrals of Sec.~\ref{sec:binding} from the
present expression.  The procedure is to expand
Eq.~\eqref{eq:3.34} in powers of $\kappa$, where $\kappa$ is the
coefficient of the nearest-neighbor coupling terms
$\bar\psi_{n}\gamma_{\mu}\psi_{n+\hat\mu}e^{i\thl{\mu\nu}}$ etc.  This
nearest-neighbor coupling term can be represented diagrammatically by a
line from the site $n$ to the site $n+\hat\mu$.  The expansion is best
described by studying an example of a term from the expansion of the
numerator of Eq.~\eqref{eq:3.34}, which will now be discussed.  An
example of a term in the expansion is represented diagrammatically in
Fig.~\ref{fig:9}.  The expression for this term is
\begin{equation}
\begin{split}
D = {} & (-1)\kappa^{4}\,
    e^{i[\thl{\mu\nu}(0,0)+\thl{\mu\nu}(1,0)+\thl{\mu\nu}(0,1)+\thl{\mu\nu}(1,1)]}\\
      & \times \bigl(\bar\psi_{00}\gamma_{\mu}\psi_{10}\bigr)
        \bigl(\bar\psi_{10}\gamma_{\nu}\psi_{11}\bigr)
        \bigl(\bar\psi_{11}\gamma_{\mu}\psi_{01}\bigr)
        \bigl(\bar\psi_{01}\gamma_{\nu}\psi_{00}\bigr),
\end{split}
\label{eq:3.36}
\end{equation}
($D$ for diagram), where the four lattice sites involved are
$(n_{0},n_{1}) = (0,0)$, $(1,0)$, $(0,1)$, and $(1,1)$; the values of
$n_{2}$ and $n_{3}$ are constant and have been suppressed in the notation
of Eq.~\eqref{eq:3.36}.  The action $A_{0}$ omits the $\kappa$ term, and
is
\begin{equation}
A_{0} = -\epsilon\sum_{n}\bar\psi_{n}\psi_{n}
  + \frac{1}{2g^{2}}\sum_{n,\mu\nu}
    \bigl(e^{if_{\mu\nu}(n)}+e^{-if_{\mu\nu}(n)}\bigr).
\label{eq:3.37}
\end{equation}

\begin{figure}[t]
  \centering
  \includegraphics[width=0.3\textwidth]{fig9}
  \refstepcounter{figure}%
  \caption*{Fig. 9. Elementary square on the lattice.}
  \label{fig:9}
\end{figure}

The calculation of $D$ has two parts, one being the integration over all
gauge fields $\thl{\mu\nu}$, the other being the calculation of the
$\psi,\bar\psi$ bracket.  These are independent calculations, i.e., $D$
factors into $D_{\theta}D_{\psi}$.  The quantity $D_{\theta}$ is
\begin{equation}
D_{\theta} = \prod_{n,\mu\nu}\int\ddth{\mu\nu}
  \exp\!\Bigl[i\bigl(\thl{\mu\nu}(0,0)+\thl{\mu\nu}(1,0)
    -\thl{\mu\nu}(1,1)-\thl{\mu\nu}(0,1)\bigr)+A_{0}\Bigr].
\label{eq:3.38}
\end{equation}

This is an example of a gauge-field average of the exponential of a line
integral over a closed loop, the loop being the loop of Fig.~\ref{fig:9}.
The $\psi$ bracket calculation can be factorized further into separate
bracket calculations for each lattice site, since $A_{0}$ contains no
terms involving $\psi$ and $\bar\psi$ coupling different lattice sites.
Consider only the four lattice sites on the loop, for simplicity.  By
moving the $\psi$'s around some (using the anticommuting rule) the
bracket becomes
\[
\begin{split}
&(\bar\psi_{00}e^{-i\thl{\mu\nu}(0,0)}\psi_{00})
 (\bar\psi_{10}e^{-i\thl{\mu\nu}(1,0)}\psi_{10})\\
&\qquad\times
 (\bar\psi_{11}e^{-i\thl{\mu\nu}(1,1)}\psi_{11})
 (\bar\psi_{01}e^{-i\thl{\mu\nu}(0,1)}\psi_{01}).
\end{split}
\]
To make a product of all possible $\psi$'s and $\bar\psi$'s means one
must have products of all possible $\psi_{00}$'s and $\bar\psi_{00}$'s,
all possible $\psi_{10}$'s and $\bar\psi_{10}$'s, etc.  In summary the
complete bracket may be written as a product of four separate brackets.
Define
\begin{equation}
D_{\psi,0} = \bigl(\bar\psi_{00}e^{i\thl{\mu\nu}(0,0)}\psi_{10}\bigr),
\label{eq:3.39}
\end{equation}
\begin{equation}
D_{\theta,0} = \bigl(\bar\psi_{10}\psi_{10}\bigr)_{0}.
\label{eq:3.40}
\end{equation}
Both $D_{\theta}$ and $D_{\psi}$ are matrices in spin space due to the
spinor indices implied for $\psi_{10}$, $\bar\psi_{10}$, $\psi_{00}$,
$\bar\psi_{00}$, and the $\gamma$ matrices.  The full bracket is simply
\begin{equation}
D_{\psi} = -\kappa^{4}\Tr\bigl(D_{\psi,0}\gamma_{\mu}
  D_{\psi,0}\gamma_{\nu}D_{\psi,0}\gamma_{\mu}D_{\psi,0}\gamma_{\nu}\bigr).
\label{eq:3.41}
\end{equation}

The matrices $D_{\psi,0}$ and $D_{\theta,0}$ are easily determined.  For
example, $D_{\theta,0}$ explicitly is
\begin{equation}
D_{\theta,0} = \delta_{\alpha\beta}
  - \epsilon\bigl(\gamma_{0}\gamma_{0}\bigr)_{\alpha\beta}
  - \epsilon\bigl(\gamma_{1}\gamma_{1}\bigr)_{\alpha\beta}
  - \epsilon\bigl(\gamma_{2}\gamma_{2}\bigr)_{\alpha\beta}
  - \epsilon\bigl(\gamma_{3}\gamma_{3}\bigr)_{\alpha\beta}.
\label{eq:3.42}
\end{equation}
The exponential can be expanded in powers of $\epsilon$; assuming the
spinors have four components only the $\epsilon^{4}$ term can produce a
product of all four $\psi_{10}$'s times all four $\bar\psi_{10}$'s; the
result is
\begin{equation}
D_{\theta,0} = \epsilon^{4}\bigl(\gamma_{0}\gamma_{1}\gamma_{2}\gamma_{3}\bigr)_{\alpha\beta},
\label{eq:3.43}
\end{equation}
(the minus sign comes from the convention that the bracket is positive
when $\bar\psi_{10}$ appears to the left of $\psi_{10}$).  A similar
calculation gives
\begin{equation}
D_{\psi,0} = \bigl(\psi_{10}\bar\psi_{10}\bigr)_{0},
\label{eq:3.44}
\end{equation}

The results of this example are easily generalized.  A term of general
order $\kappa^{m}$ is nonzero only if the nearest-neighbor couplings
combine to form closed loops (the lattice site at the endpoint of an open
line would have an extra $\psi_{n}$ or $\bar\psi_{n}$ so the bracket at
$n$ would give 0).  The bracket calculation for a closed loop gives a
trace involving $\kappa$ times a $\gamma$ matrix for each line in the
loop and $D$'s for each lattice site in the loop (except the points $n$
and $0$ where there are currents).  The average over gauge fields
involves an exponential of a sum of $\thl{\mu\nu}$'s around each loop.
There can be any number of loops.
